\documentclass[11pt,a4paper]{article}

\usepackage[utf8]{inputenc}
\usepackage[T1]{fontenc}
\usepackage{textcomp}
\usepackage{lmodern}
\usepackage[french]{babel}
\usepackage[margin=2.2cm,top=2.4cm,bottom=2.2cm]{geometry}
\usepackage{microtype}
\usepackage{xcolor}
\usepackage{enumitem}
\usepackage{titlesec}
\usepackage{fancyhdr}
\usepackage{tikz}
\usepackage{pgfplots}
\usepackage[most]{tcolorbox}
\usepackage[hidelinks]{hyperref}
\pgfplotsset{compat=1.16}
\usetikzlibrary{arrows.meta,positioning,shapes.geometric,fit,calc,backgrounds}

% ---------- palette ----------
\definecolor{accent}{HTML}{0B5D8A}
\definecolor{accentlight}{HTML}{E8F1F7}
\definecolor{good}{HTML}{2E7D32}
\definecolor{goodlight}{HTML}{E6F4E7}
\definecolor{bad}{HTML}{B3261E}
\definecolor{badlight}{HTML}{FBE9E7}
\definecolor{warn}{HTML}{B45309}
\definecolor{warnlight}{HTML}{FFF3E0}
\definecolor{ink}{HTML}{1F2933}
\definecolor{mute}{HTML}{6B7280}

% ---------- headings ----------
\titleformat{\section}{\sffamily\Large\bfseries\color{accent}}{\thesection}{0.7em}{}[{\color{accent}\titlerule[0.8pt]}]
\titlespacing*{\section}{0pt}{2.6ex plus 0.5ex}{1.4ex}
\titleformat{\subsection}{\sffamily\large\bfseries\color{ink}}{}{0em}{}
\titlespacing*{\subsection}{0pt}{2ex}{0.8ex}

\pagestyle{fancy}\fancyhf{}
\fancyhead[L]{\sffamily\footnotesize\bfseries\color{accent} PER MEA S9P-2026}
\fancyhead[R]{\sffamily\footnotesize\color{mute} Rapport d'avancement --- 30 septembre 2026}
\fancyfoot[C]{\sffamily\footnotesize\color{mute}\thepage}
\renewcommand{\headrule}{\color{accent}\hrule width\headwidth height0.6pt}
\setlength{\parskip}{0.6em}\setlength{\parindent}{0pt}
\renewcommand{\arraystretch}{1.2}

% ---------- boxes ----------
\tcbset{boxstyle/.style={enhanced,arc=2.5pt,boxrule=0pt,leftrule=3.5pt,left=9pt,right=9pt,top=6pt,bottom=6pt,breakable}}
\newtcolorbox{info}[1][]{boxstyle,colback=accentlight,colframe=accent,#1}
\newtcolorbox{alert}[1][]{boxstyle,colback=badlight,colframe=bad,#1}
\newtcolorbox{ok}[1][]{boxstyle,colback=goodlight,colframe=good,#1}
\newtcolorbox{caution}[1][]{boxstyle,colback=warnlight,colframe=warn,#1}

% KPI tile: \kpi{value}{label}{color}
\newcommand{\kpi}[3]{%
\begin{tcolorbox}[enhanced,colback=#3!8,colframe=#3,boxrule=0.6pt,arc=3pt,halign=center,
  left=2pt,right=2pt,top=5pt,bottom=5pt]
{\sffamily\bfseries\LARGE\color{#3}#1}\\[-1pt]{\sffamily\footnotesize\color{ink}#2}
\end{tcolorbox}}

% Before / after pair: \fix{title}{before}{after}
\newcommand{\fix}[3]{%
\par\medskip
{\sffamily\bfseries\color{ink}#1}\par\vspace{2pt}
\noindent\begin{minipage}[t]{0.455\linewidth}
\begin{tcolorbox}[enhanced,colback=badlight,colframe=bad,boxrule=0pt,leftrule=3pt,arc=2pt,
 left=7pt,right=6pt,top=4pt,bottom=4pt,height=2.6cm,valign=center]
{\sffamily\scriptsize\bfseries\color{bad} AVANT}\\[1pt]{\small #2}
\end{tcolorbox}\end{minipage}%
\hfill\raisebox{1.1cm}{\tikz\draw[-{Stealth[length=2.6mm]},line width=1.3pt,color=mute](0,0)--(0.45,0);}\hfill
\begin{minipage}[t]{0.455\linewidth}
\begin{tcolorbox}[enhanced,colback=goodlight,colframe=good,boxrule=0pt,leftrule=3pt,arc=2pt,
 left=7pt,right=6pt,top=4pt,bottom=4pt,height=2.6cm,valign=center]
{\sffamily\scriptsize\bfseries\color{good} APRÈS}\\[1pt]{\small #3}
\end{tcolorbox}\end{minipage}\par\medskip}

% Process block for TikZ diagrams
\tikzset{
  blk/.style={rounded corners=3pt,draw=accent,fill=accentlight,thick,align=center,
    minimum height=1.05cm,text width=2.35cm,font=\sffamily\footnotesize,inner sep=3pt},
  blkbad/.style={blk,draw=bad,fill=badlight},
  blkgood/.style={blk,draw=good,fill=goodlight},
  arr/.style={-{Stealth[length=2.4mm]},thick,color=mute},
  tag/.style={circle,fill=bad,text=white,font=\sffamily\bfseries\scriptsize,inner sep=1.6pt},
  tagok/.style={circle,fill=good,text=white,font=\sffamily\bfseries\scriptsize,inner sep=1.6pt},
}

\pgfplotsset{
  barstyle/.style={ybar,bar width=14pt,width=\linewidth,height=5.2cm,ymin=0,
    axis lines*=left,ymajorgrids,grid style={color=black!10},tick style={draw=none},
    every axis label/.style={font=\sffamily\footnotesize},
    tick label style={font=\sffamily\footnotesize},
    nodes near coords,every node near coord/.append style={font=\sffamily\scriptsize},
    enlarge x limits=0.25,
  }
}

\begin{document}

% ================= TITLE =================
\begin{tcolorbox}[enhanced,colback=accent,colframe=accent,coltext=white,arc=3pt,
  left=16pt,right=16pt,top=14pt,bottom=14pt,boxrule=0pt]
{\sffamily\footnotesize\bfseries\color{white!80!accent} RAPPORT D'AVANCEMENT \quad$\cdot$\quad 30 SEPTEMBRE 2026}\\[0.5em]
{\sffamily\Huge\bfseries\color{white} Fiabiliser le drone voilier}\\[0.15em]
{\sffamily\large\color{white} Audit, corrections et simulateur --- PER MEA S9P-2026}\\[0.6em]
{\sffamily\small\color{white} Haitam Kilani \quad$\cdot$\quad partie embarquée}
\end{tcolorbox}

\vspace{0.3em}
\begin{info}
\textbf{En bref.} Au départ, tous les tests passaient, mais ils ne vérifiaient pas grand-chose. J'ai cherché ce qu'ils ne
contrôlaient pas, reproduit chaque défaut avant de le corriger, et étendu le simulateur pour qu'il représente enfin les
perturbations réelles. Un point reste bloquant avant tout essai autonome : la course réelle du safran, qui ne se règle
pas dans le code mais se mesure au banc. Rien de ce travail n'est encore committé.
\end{info}

% ================= 1 =================
\section{Ma méthode}

Je n'ai retenu un défaut que lorsque j'ai pu le montrer en l'exécutant. Chaque constat suit le même parcours.

\vspace{0.3em}
\begin{center}
\begin{tikzpicture}[node distance=7mm]
\node[blk] (a) {\textbf{1. Relire}\\le code et les tests};
\node[blk,right=of a] (b) {\textbf{2. Reproduire}\\le défaut par un test};
\node[blk,right=of b] (c) {\textbf{3. Corriger}\\sans toucher au matériel};
\node[blkgood,right=of c] (d) {\textbf{4. Vérifier}\\tests et re-mutation};
\draw[arr](a)--(b);\draw[arr](b)--(c);\draw[arr](c)--(d);
\end{tikzpicture}
\end{center}

Le travail couvre cinq zones du projet, résumées ci-dessous. Le détail suit dans les sections suivantes.

\vspace{0.3em}
\begin{center}
\begin{tikzpicture}
\node[blk,text width=3.05cm,minimum height=2.3cm,anchor=north west] (f) at (0,0)
 {\textbf{Firmware}\\[2pt]\scriptsize navigation,\\radio LoRa,\\gestion des modes};
\node[blk,text width=3.05cm,minimum height=2.3cm,anchor=north west] (s) at (3.3,0)
 {\textbf{Simulateur}\\[2pt]\scriptsize courant, dérive,\\vent réel,\\scénarios isolés};
\node[blk,text width=3.05cm,minimum height=2.3cm,anchor=north west] (h) at (6.6,0)
 {\textbf{Harnais adverse}\\[2pt]\scriptsize chiffres refaits,\\statistiques\\corrigées};
\node[blk,text width=3.05cm,minimum height=2.3cm,anchor=north west] (i) at (9.9,0)
 {\textbf{Interface (IHM)}\\[2pt]\scriptsize trames, ordre des\\commandes, base\\de données};
\node[blk,text width=4.7cm,minimum height=1.05cm,anchor=north west] (t) at (0,-2.55)
 {\textbf{Tests} : suites renforcées et remises en marche};
\node[blk,text width=7.9cm,minimum height=1.05cm,anchor=north west] (d) at (5.05,-2.55)
 {\textbf{Documentation} : écarts entre la description du projet et le code corrigés};
\end{tikzpicture}
\end{center}

% ================= 2 =================
\section{Point de départ : la chaîne GPS}

Le 23 septembre, j'avais montré que les problèmes de terrain venaient du GPS et qu'ils existaient dans
\emph{toutes} les versions du firmware. Le bateau n'a pas de compas : à basse vitesse, son seul cap est la route mesurée
par le GPS, et le programme traitait comme valides des données périmées.

\begin{center}
\begin{tikzpicture}[node distance=6mm]
\node[blk,text width=2.2cm] (g) {\textbf{Récepteur GPS}};
\node[blk,text width=2.2cm,right=of g] (p) {\textbf{Pilote}\\validité du fix};
\node[blk,text width=2.2cm,right=of p] (c) {\textbf{Calcul du cap}\\et de la vitesse};
\node[blk,text width=2.2cm,right=of c] (n) {\textbf{Navigation}\\et trim du safran};
\node[blk,text width=1.8cm,right=of n] (a) {\textbf{Safran,\\voile}};
\draw[arr](g)--(p);\draw[arr](p)--(c);\draw[arr](c)--(n);\draw[arr](n)--(a);
\node[tag,above=1mm of g.north east]{1};
\node[tag,above=1mm of p.north east]{2};
\node[tag,above=1mm of c.north east]{3};
\node[tag,above=1mm of n.north east]{4};
\node[text width=2.2cm,align=center,font=\sffamily\scriptsize\color{bad},below=4mm of g]{Réglage du récepteur inadapté};
\node[text width=2.2cm,align=center,font=\sffamily\scriptsize\color{bad},below=4mm of p]{Perte de fix jamais détectée};
\node[text width=2.4cm,align=center,font=\sffamily\scriptsize\color{bad},below=4mm of c]{Cap à 0° au démarrage, valeurs figées};
\node[text width=2.4cm,align=center,font=\sffamily\scriptsize\color{bad},below=4mm of n]{Pilote sur un faux cap};
\end{tikzpicture}
\end{center}

\begin{ok}
\textbf{Ce que j'ai mis en place.} Un fix n'est valide que s'il a moins de 2 secondes, au moins 5 satellites et une
précision suffisante. Le cap n'est accepté qu'après confirmation sur 3 mesures, avec une marge de vitesse pour éviter
qu'il clignote. Si le fix tombe, le bateau se met en neutre et repart seul ; si le cap est inconnu, il avance en ligne
droite à l'hélice. J'ai aussi ajouté un réglage automatique du décalage du safran, plus un chien de garde qui détecte
une manœuvre bloquée.
\end{ok}

\subsection{Effet du réglage automatique du safran}

J'ai simulé 854 cas d'un décalage mécanique du safran, de $-30^\circ$ à $+30^\circ$. Au-delà de $20^\circ$, le bateau
tournait en rond ; c'est maintenant réglé dans 99,4~\% des cas.

\begin{center}
\begin{tikzpicture}
\begin{axis}[barstyle,ylabel={Missions réussies (\%)},symbolic x coords={{Décalage $\leq 4^\circ$},{5 à $19^\circ$},{$\geq 20^\circ$}},
 xtick=data,ymax=118,legend style={font=\sffamily\scriptsize,at={(0.5,1.02)},anchor=south,legend columns=2,draw=none},
 nodes near coords={\pgfmathprintnumber[fixed,precision=1]{\pgfplotspointmeta}},height=5.4cm,width=0.85\linewidth,bar width=18pt]
\addplot[fill=bad!70,draw=none] coordinates {({Décalage $\leq 4^\circ$},100) ({5 à $19^\circ$},87.4) ({$\geq 20^\circ$},9.7)};
\addplot[fill=good!80,draw=none] coordinates {({Décalage $\leq 4^\circ$},100) ({5 à $19^\circ$},100) ({$\geq 20^\circ$},99.4)};
\legend{Avant,Après}
\end{axis}
\end{tikzpicture}
\end{center}

Un couloir de navigation ramené de 100 à 30~m réduit l'écart maximal à la route de 103 à 57~m, pour 4~\% de durée en plus.

% ================= 3 =================
\section{Corrections dans le firmware}

\fix{Un angle extrême bloquait le calcul}
{La normalisation d'un angle tournait en boucle : environ 6~millions de tours pour une valeur extrême, soit 14,6~ms sur mon PC et probablement plusieurs secondes sur la carte.}
{Le calcul est direct et de durée constante. Deux fonctions voisines corrigeaient l'angle une seule fois ; elles sont aussi réparées (avant : 9\,639$^\circ$ pour un vent de 9\,999$^\circ$).}

\fix{Des données radio prises pour argent comptant}
{Le vent reçu n'était pas contrôlé, et une trame tronquée donnait un waypoint en (0,0), très loin de la zone d'essai, vers lequel le bateau naviguait sans erreur.}
{Le vent est vérifié et ramené entre 0 et 360$^\circ$. Les waypoints sont contrôlés (latitude, longitude, rejet de (0,0)) et la mission est acceptée en entier ou refusée en entier.}

\fix{L'état de navigation survivait au mode Manuel}
{Après un passage en Manuel, les repères du couloir et la phase d'anti-empannage restaient en mémoire.}
{Ils sont remis à zéro en quittant le mode Automatique. Le décalage appris du safran est conservé, car il dépend du matériel.}

\fix{Un plafond de puissance trop haut}
{Le régime automatique pouvait monter à 2000\,\textmu s (pleine puissance).}
{Il est plafonné à 1850\,\textmu s, la limite prévue pour l'autonomie.}

% ================= 4 =================
\section{Le point ouvert : la course du safran}

En relisant le code, j'ai vu que la limite de $\pm 20^\circ$ que j'avais validée au banc le 4~juin, puis installée sur le
bateau, n'existe plus. Le mode automatique autorise maintenant $\pm 110^\circ$.

\begin{center}
\begin{tikzpicture}[x=0.5cm]
\draw[thick,color=mute](-11,0)--(11,0);
\foreach \x in {-11,-9,-2,0,2,9,11}{\draw[color=mute](\x,-0.08)--(\x,0.08);}
\draw[line width=7pt,color=good,opacity=0.85](-2,0.6)--(2,0.6);
\draw[line width=7pt,color=warn,opacity=0.85](-9,1.25)--(9,1.25);
\draw[line width=7pt,color=bad,opacity=0.85](-11,1.9)--(11,1.9);
\node[font=\sffamily\scriptsize,anchor=west] at (2.4,0.6){$\pm 20^\circ$ : validé au banc (4 juin)};
\node[font=\sffamily\scriptsize,anchor=west] at (9.4,1.25){$\pm 90^\circ$ : course notée};
\node[font=\sffamily\scriptsize,anchor=south] at (0,2.3){$\pm 110^\circ$ : code actuel};
\node[font=\sffamily\scriptsize,anchor=north,color=mute] at (-9,-0.15){$-90^\circ$};
\node[font=\sffamily\scriptsize,anchor=north,color=mute] at (9,-0.15){$+90^\circ$};
\node[font=\sffamily\scriptsize,anchor=north,color=mute] at (0,-0.15){$0^\circ$};
\end{tikzpicture}
\end{center}

J'ai balayé 5\,184 géométries de navigation. Avec un réglage neutre, la commande ne dépasse jamais $90^\circ$
(pire cas : $80,1^\circ$). Le danger apparaît quand le décalage appris est chargé, ce qui arrive naturellement, car il
persiste d'un waypoint à l'autre.

\begin{center}
\begin{tikzpicture}
\begin{axis}[barstyle,ylabel={Commandes (\%)},symbolic x coords={{$0^\circ$},{$+20^\circ$},{$+40^\circ$},{$-40^\circ$}},
 xtick=data,xlabel={Décalage appris au départ},ymax=17,
 legend style={font=\sffamily\scriptsize,at={(0.5,1.02)},anchor=south,legend columns=2,draw=none},
 nodes near coords={\pgfmathprintnumber[fixed,precision=1]{\pgfplotspointmeta}},height=5.2cm,width=0.85\linewidth,bar width=15pt]
\addplot[fill=warn!80,draw=none] coordinates {({$0^\circ$},0) ({$+20^\circ$},0) ({$+40^\circ$},10.4) ({$-40^\circ$},13.2)};
\addplot[fill=bad!80,draw=none] coordinates {({$0^\circ$},0) ({$+20^\circ$},0) ({$+40^\circ$},0.7) ({$-40^\circ$},10.4)};
\legend{Dépasse $90^\circ$,Bloqué à $110^\circ$}
\end{axis}
\end{tikzpicture}
\end{center}

\begin{alert}
\textbf{Décision.} Je n'ai modifié aucune constante d'actionneur. Savoir si $90^\circ$ ou $110^\circ$ est une butée
mécanique du winch ne se tranche pas en logiciel : il faut le mesurer au banc. J'ai en revanche corrigé la documentation,
qui annonçait encore $\pm 20^\circ$ à trois endroits.
\end{alert}

% ================= 5 =================
\section{Le simulateur}

Le simulateur donnait la même valeur au cap du bateau et à sa route mesurée par le GPS. Or c'est justement l'écart
entre les deux, causé par le courant et la dérive, que le bateau ne peut pas détecter sans compas. « Six scénarios sur
six réussissent » ne disait donc presque rien du risque principal.

\begin{center}
\begin{tikzpicture}
\begin{scope}
\node[font=\sffamily\bfseries\footnotesize] at (1.9,2.7){Avant : cap = route};
\draw[-{Stealth},line width=2pt,color=accent](0,0)--(3.6,1.8);
\node[font=\sffamily\scriptsize,anchor=south east] at (3.7,1.9){cap et route};
\end{scope}
\begin{scope}[xshift=7cm]
\node[font=\sffamily\bfseries\footnotesize] at (2.2,2.7){Après : cap $\neq$ route};
\draw[-{Stealth},line width=2pt,color=accent](0,0)--(2.6,1.3);
\draw[-{Stealth},line width=2pt,color=warn](2.6,1.3)--(4.3,0.5);
\draw[-{Stealth},line width=2.4pt,color=bad,dashed](0,0)--(4.3,0.5);
\node[font=\sffamily\scriptsize,color=accent,anchor=south east] at (2.4,1.4){cap};
\node[font=\sffamily\scriptsize,color=warn,anchor=south west] at (3.4,1.1){courant, dérive};
\node[font=\sffamily\scriptsize,color=bad,anchor=north] at (2.4,0.05){route mesurée par le GPS};
\end{scope}
\end{tikzpicture}
\end{center}

J'ai ajouté le courant et la dérive sous le vent, un vent imposé au choix et une option qui va chercher la prévision de
vent réelle. Sans option, les résultats sont identiques à avant : le scénario~6 donne toujours 688~s. Le simulateur
renvoie maintenant un résultat de réussite ou d'échec, ce qui en fait un test de non-régression.

\noindent
\begin{minipage}[t]{0.485\linewidth}
\begin{info}\textbf{Corrections d'isolement}\\\small Les six scénarios partageaient un état de navigation et
n'étaient pas indépendants. Chacun a maintenant le sien. Le choix d'un scénario par son numéro fonctionne, et les
fichiers de compilation sont à jour.\end{info}
\end{minipage}\hfill
\begin{minipage}[t]{0.485\linewidth}
\begin{info}\textbf{Nouvelles options}\\\small Courant (vitesse et direction), dérive, vent imposé (direction et
vitesse) et prévision réelle (Windy).\end{info}
\end{minipage}

\begin{caution}
\textbf{Windy : pas encore opérationnel.} La prévision est bien demandée, mais Windy répond « clé invalide ». Il faut
sans doute une clé de type \emph{Point Forecast}. La conversion des données de vent a été vérifiée avec une fausse
réponse.
\end{caution}

% ================= 6 =================
\section{Le harnais adverse}

Ce banc fait tourner la vraie navigation avec des perturbations que le simulateur ignorait. J'ai constaté que les
chiffres publiés ne se reproduisaient pas, avec 9 à 13 points d'écart sur du code déterministe, et que le tirage
aléatoire changeait à chaque valeur testée. Après correction, j'ai refait les 200 missions par point et recalculé.

\begin{center}
\begin{minipage}{0.48\linewidth}
\begin{tikzpicture}
\begin{axis}[barstyle,ylabel={Missions réussies (\%)},symbolic x coords={{$0$},{$0,25$ m/s}},xtick=data,
 xlabel={Courant},ymax=92,height=5cm,width=\linewidth,bar width=22pt]
\addplot[fill=accent!80,draw=none] coordinates {({$0$},73.5) ({$0,25$ m/s},25.5)};
\end{axis}
\end{tikzpicture}
\end{minipage}\hfill
\begin{minipage}{0.48\linewidth}
\begin{tikzpicture}
\begin{axis}[barstyle,symbolic x coords={{$0^\circ$},{$10^\circ$},{$20^\circ$}},xtick=data,
 xlabel={Erreur sur la direction du vent},ymax=92,height=5cm,width=\linewidth,bar width=18pt]
\addplot[fill=accent!80,draw=none] coordinates {({$0^\circ$},75.5) ({$10^\circ$},67.5) ({$20^\circ$},49.0)};
\end{axis}
\end{tikzpicture}
\end{minipage}
\end{center}

\noindent
\begin{minipage}[t]{0.32\linewidth}
\begin{ok}\sffamily\small\textbf{Solide}\\Le courant et l'erreur de vent font chuter la convergence.\end{ok}
\end{minipage}\hfill
\begin{minipage}[t]{0.32\linewidth}
\begin{info}\sffamily\small\textbf{Sans effet}\\La vitesse du winch ne change rien de mesurable.\end{info}
\end{minipage}\hfill
\begin{minipage}[t]{0.32\linewidth}
\begin{caution}\sffamily\small\textbf{Non concluant}\\Le bruit GPS : l'ancien balayage ne permet pas de conclure.\end{caution}
\end{minipage}

J'ai aussi corrigé un plantage lorsque le fichier de résultats n'est pas inscriptible, et régénéré le document de résultats
avec le nombre de missions et les intervalles de confiance.

% ================= 7 =================
\section{L'interface (IHM)}

Chaque maillon de la chaîne entre le bateau et l'écran de l'opérateur avait un défaut. Les pastilles indiquent où j'ai
agi.

\begin{center}
\begin{tikzpicture}[node distance=6mm]
\node[blk,text width=1.9cm] (b) {\textbf{Bateau}};
\node[blk,text width=1.9cm,right=of b] (t) {\textbf{Émetteur} LoRa};
\node[blk,text width=2.2cm,right=of t] (s) {\textbf{Liaison série}};
\node[blk,text width=2.2cm,right=of s] (m) {\textbf{Base de données}};
\node[blk,text width=1.9cm,right=of m] (w) {\textbf{Carte web}};
\draw[arr](b)--(t);\draw[arr](t)--(s);\draw[arr](s)--(m);\draw[arr](m)--(w);
\node[tagok,above=1mm of s.north east]{1};
\node[tagok,above=1mm of m.north east]{2};
\node[tagok,above=1mm of w.north east]{3};
\end{tikzpicture}
\end{center}

\noindent
\begin{minipage}[t]{0.32\linewidth}
\begin{ok}\sffamily\small\textbf{1. Trames perdues}\\Une trame coupée entre deux lectures était rejetée en entier. Elle est maintenant reconstituée.\end{ok}
\end{minipage}\hfill
\begin{minipage}[t]{0.32\linewidth}
\begin{ok}\sffamily\small\textbf{2. Ordre et accès}\\Les commandes partent dans l'ordre d'arrivée (un Stop ne perd plus contre un Navigate). La base n'est plus ouverte à tout le réseau.\end{ok}
\end{minipage}\hfill
\begin{minipage}[t]{0.32\linewidth}
\begin{ok}\sffamily\small\textbf{3. Vérifications}\\Une route avec un point invalide est refusée en entier. La prévision de vent est affichée à l'opérateur, jamais envoyée toute seule.\end{ok}
\end{minipage}

% ================= 8 =================
\section{Les tests}

\noindent
\begin{minipage}{0.24\linewidth}\kpi{4/4}{tests du firmware}{good}\end{minipage}\hfill
\begin{minipage}{0.24\linewidth}\kpi{61}{tests de l'interface}{good}\end{minipage}\hfill
\begin{minipage}{0.24\linewidth}\kpi{20}{tests du banc GPS}{good}\end{minipage}\hfill
\begin{minipage}{0.24\linewidth}\kpi{6/6}{scénarios de simulation}{good}\end{minipage}

\vspace{0.4em}
Ces résultats étaient déjà verts avant mon travail, ce qui prouve seulement qu'ils étaient trop faibles. Pour le
montrer, j'ai injecté volontairement une erreur (test de mutation) : en multipliant par 28 le gain de la propulsion, aucune
suite ne réagissait.

\begin{center}
\begin{tikzpicture}[node distance=8mm]
\node[blk,text width=2.6cm] (e) {\textbf{Erreur injectée}\\gain $\times$ 28};
\node[blkbad,text width=3cm,right=of e] (av) {\textbf{Avant}\\4 suites vertes, l'erreur passe};
\node[blkgood,text width=3cm,right=of av] (ap) {\textbf{Après}\\le test devient rouge};
\draw[arr](e)--(av);\draw[arr](av)--(ap);
\node[font=\sffamily\scriptsize,color=mute] at ($(av.east)+(0,0.85)$) {test renforcé};
\end{tikzpicture}
\end{center}

J'ai remplacé les vérifications trop faibles par des valeurs exactes et ajouté des tests sur les limites des angles et
sur les 16~waypoints. La suite du banc GPS ne s'exécutait pas du tout (une erreur de structure l'arrêtait avant le premier
test) : elle tourne maintenant.

\begin{alert}
\textbf{Encore sans aucun test :} le choix du mode de fonctionnement, l'écriture vers les servos et l'ESC, et la
réception des commandes par radio.
\end{alert}

% ================= 9 =================
\section{Ce qui reste à faire}

\begin{center}
\begin{tikzpicture}
\def\row#1#2#3#4{\node[draw=#4,fill=#4!8,rounded corners=2pt,thick,text width=14.4cm,inner sep=5pt,align=left,font=\small,anchor=north west] at (0,#1) {\textbf{#2}\ \ #3};}
\row{0}{Prioritaire}{Mesurer au banc la course réelle du winch et trancher entre $\pm 20^\circ$, $\pm 90^\circ$ et $\pm 110^\circ$.}{bad}
\row{-1.0}{À faire}{Obtenir une clé Windy valide et lancer le simulateur avec la prévision réelle.}{warn}
\row{-2.0}{À faire}{Écrire des tests pour le choix du mode, l'écriture vers les servos et la réception radio.}{warn}
\row{-3.0}{À faire}{Fixer l'erreur de vent tolérable et la vérifier en conditions réelles.}{warn}
\row{-4.0}{À faire}{Mesurer sur la carte le temps de blocage de l'ancienne boucle d'angle (je ne l'ai qu'estimé).}{warn}
\row{-5.0}{À faire}{Examiner les deux échecs restants du balayage de décalage du safran, puis committer tout ce travail.}{warn}
\end{tikzpicture}
\end{center}

\begin{info}
\textbf{Risque accepté pour l'instant.} La radio et l'interface web n'exigent pas d'identification : toute commande bien
formée est exécutée. Pour un projet de terrain aux essais encadrés, je le laisse de côté ; passer les commandes en
requêtes protégées par un secret partagé est prévu plus tard.
\end{info}

\end{document}
